\section{LDP Path Counting}
\label{sec:k-lp}

 
Now, we extend our LDP solution for walk counting to support path counting.  
As discussed in Section~\ref{subsec:prelim_pattern_counting}, path counting imposes an additional constraint compared with walk counting: each node must appear at most once along the path. 
Even in the non-private, centralized setting, this constraint introduces significant challenges.
As shown in the prior work~\cite{jokic2022number}, in theory, walk counting can be solved asymptotically more efficiently than path counting.
Moving towards the LDP setting, the problem becomes even more challenging, as enforcing the uniqueness constraint for a path would inherently require information sharing across multiple nodes, which is fundamentally difficult to achieve under the LDP model.


To address this issue, in this section, we introduce a protocol named random marking, which 
can be seamlessly integrated into our walk counting mechanisms to enable private $k$-line path counts estimation under LDP. 
The resulting mechanism achieves a utility bound of $\tilde{O}(\sqrt{N}d(G)^{k})$, representing only a multiplicative overhead of $\tilde{O}(d(G))$ compared to the corresponding $k$-line walk counting utility of $\tilde{O}(\sqrt{N}d(G)^{k-1})$.



\subsection{Random Marking: Transforming Path Counting into Walk Counting 
}
\label{subsec:k-lp_randommarking}


% \red{As mentioned, detecting node repetition in a path under the LDP setting is inherently difficult due to the lack of global visibility. To address this challenge, we adopt a different strategy: instead of detecting repetitions after the fact, we proactively prevent node repetition from the outset. Our solution is to impose a stricter constraint. Specifically, for each node, we not only require that it appear at most once in a walk, but also that it must be in a predefined position within the walk. }


As highlighted above, detecting node repetition in a path under the LDP setting is inherently difficult due to limited global visibility. 
To address this challenge, {we draw inspiration from color-coding~\cite{alon1995color}, which colors nodes and searches for patterns with distinct node colors}, and adopt a different strategy: rather than detecting repetitions retrospectively, we proactively prevent them by imposing a stricter constraint from the outset.
Specifically, for each node, we require that it not only appear at most once in a walk, but also occupy a predefined position within it. 
{Interestingly, this idea was originally introduced to improve runtime for non-private path and cycle detection, whereas we leverage it here to address the challenges in the LDP setting, where only local views are available.}




To illustrate the high-level idea, we first consider the non-private setting and defer the DP-relevant operations to the next subsection.
We begin with an extra marking round, in which each node $v_i$ samples a random mark $r_i$ uniformly from $\{0, 1, \dots, k\}$ and sends it to its neighbors. 
This mark enforces that the node may appear only as the $(r_i + 1)$-th node in any $k$-line path being counted. 
For example, if a node $v_j$ draws $r_j = 0$, it can only serve as the starting node of a valid $k$-line path. 
We then apply the $(k-1)$-round walk counting mechanism $\mathcal{M}_{\mathrm{k\text{-}wk}}$, with an added constraint: each node $v_i$ participates only in the $r_i$-th round of computation.
At all other rounds $\ell \ne r_i$, the node simply outputs $0$.
More precisely, in each round $\ell \in \{1, \dots, k-1\}$, nodes $v_i \in V$ with $r_i = \ell$ compute the number of $\ell$-line walks ending at themselves with node marks $(0, 1, \dots, r_i)$.
This is achieved by aggregating counts from active neighbors that participated in the previous round.
Additionally, nodes with mark $k - 1$ complete the last line by multiplying their count by the number of their neighbors marked with $k$. 
Notably, since $\mathcal{M}_{\mathrm{k\text{-}wk}}$ runs for only $k-1$ rounds, nodes marked with $0$ or $k$ do not participate in subsequent computation. 
To distinguish this from the DP mechanism, we denote the resulting non-private algorithm as $\widehat{\mathcal{M}}_{\mathrm{k\text{-}pt}}$, where each local client executes the function $\widehat{\mathcal{R}}^{(\ell)}_i$ at round $\ell$. 
The detailed algorithm is provided in Algorithm~\ref{alg:k-lp_nondp}.
Here, we use $\widehat{\mathcal{X}}_i^{(\ell)}$ to denote the non-private output of node $v_i$ in round $\ell$, in distinction from the DP-compliant result ${\mathcal{X}}_i^{(\ell)}$.




\begin{algorithm}
\small
\caption{Local client function $\widehat{\mathcal{R}}^{(\ell)}$ of $\widehat{\mathcal{M}}_{\mathrm{k\text{-}pt}}$}
\label{alg:k-lp_nondp}
\SetKwFunction{rnew}{$\widehat{\mathcal{R}}^{(\ell)}$}
\SetNoFillComment

\ForEach{$i \in \{1, \dots, N\}$}{
    $\widehat{\mathcal{X}}^{(0)}_i \leftarrow 1$
}


\SetKwProg{Fn}{Function}{:}{}
\Fn{\rnew{$G_i$}}{


\If{$\ell = \text{mark}$}{
    \tcc{Randomly sample a mark}
    $r_i \gets \textsf{Unif}(0,k)$

     \textbf{Send out} $r_i$ to $\mathcal{N}(i)$
   
    }

\If{$\ell = r_i$}{

Received $\widehat{\mathcal{X}}^{(\ell-1)}_j$ for each $j \in \mathcal{N}(i)$


  
        $\widehat{\mathcal{X}}^{(\ell)}_i 
        \leftarrow 
        \sum_{j \in \mathcal{N}(i)} 
        \widehat{\mathcal{X}}^{(\ell-1)}_j \, \mathbf{1}\{r_j = {\ell -1}\}$\\
    % }

    \If{$\ell = k-1$}{
        $\widehat{\mathcal{X}}^{(\ell)}_i 
        \leftarrow 
        \widehat{\mathcal{X}}^{(\ell)}_i \cdot \sum_{j \in \mathcal{N}(i)} \mathbf{1}\{r_j = k\}$\\
    }

}
\Else{
     $\widehat{\mathcal{X}}^{(\ell)}_i \leftarrow 0$
}


 \textbf{Send out} $\widehat{\mathcal{X}}^{(\ell)}_i$ to $\mathcal{N}(i)$ and $\mathcal{A}$
}
\end{algorithm}



With this design, it is straightforward to ensure that each node $v_i \in V$ appears at most once in a walk, since its contribution is limited to being the $(r_i+1)$-th node in any walk. 
% This enforces the non-repetition requirement for path counting. 
This enforces the requirement that nodes do not repeat in a path.
However, it also introduces the problem of underestimation. 
Specifically, a $k$-line path in the graph is counted only if all nodes along the path receive marks exactly from $0$ to $k$. 
Since node marks are sampled independently, the probability that a $k$-line path is sampled is $(k+1)^{-(k+1)}$. 
Consequently, at the analyzer $\mathcal{A}$, the result is rescaled by $(k+1)^{(k+1)}$:
\begin{equation*}
     \widehat{\mathcal{M}}_{\mathrm{k\text{-}pt}}(G) = (k+1)^{k+1} \, 
     \sum_{i=1}^N \widehat{\mathcal{X}}^{(k-1)}_i.
\end{equation*}
{Random marking introduces variance in  $\widehat{\mathcal{M}}_{\mathrm{k\text{-}pt}}(G)$ relative to the true $k$-line path count $P_k$. 
For utility analysis, we capture this effect as an additive sampling error, which we bound tightly below. 
Importantly, this sampling error is accounted for in the overall error of our LDP mechanism introduced in Section~\ref{subsec:k-lp_dp}.
}





\begin{lemma}
\label{lem:k-lp_sampling_bound}
For any graph $G$, path length $k$, and parameter $\beta$, 
the $k$-round non-private algorithm $\widehat{\mathcal{M}}_{\mathrm{k\text{-}pt}}$ satisfies $
    \mathbb{E}[\widehat{\mathcal{M}}_{\mathrm{k\text{-}pt}}(G)] = P_k $,
and with probability at least $1 - \beta$,
\begin{equation*}
    |\widehat{\mathcal{M}}_{\mathrm{k\text{-}pt}}(G)-P_k| \leq
    (k+1) \sqrt{N\beta^{-1}}(d(G)+k+1)^k
    = O(\sqrt{N} d(G)^k).
\end{equation*}
\end{lemma}




% \red{With this design, it becomes trivial to guarantee that each node appears at most once in a walk, as it is only permitted to appear in a predefined position. This positional constraint inherently enforces the non-repetition property required for path counting. However, it will also cause the underestimation of ... More precisely, each line path will only have a probability... Therefore, in the side of analyzer..., rescle the result by ... This will lead to an sampling error. Fortunately, below, we will show the sampling error can be well-bounded. }






\subsection{Integrating DP and Efficiency Optimization}
\label{subsec:k-lp_dp}


% With algorithm..., we successfully reduce ... into a problem of ... Then, we can utilize the idea of LDP algorithm of ... to privatize ... Here, before round ..., the analyzer will ..., and then, 
% after computing ..., the node v will add the Laplace noise to ..., But a difference here is that, since we use the random mark, one edge will at most appear in one round of computation. As a result, we do not need to divide the privacy budget ... With above idea, we have the ....



With Algorithm~\ref{alg:k-lp_nondp}, we successfully reduce the problem of $k$-line path counting to that of $k$-line walk counting. 
We can then utilize the idea of the edge-LDP mechanism $\mathcal{M}_{\mathrm{k\text{-}wk}}$ to privatize node outputs from Algorithm~\ref{alg:k-lp_nondp}.
To begin with, publishing node marks incurs no privacy loss, as the marks are independent of any sensitive data.
Before each round $\ell$, the analyzer collects the outputs from active nodes, computes the global maximum $\|\mathcal{X}^{(\ell-1)}\|_\infty$, and broadcasts it to all nodes.
Then in round $\ell$, after aggregating the result from its neighbors in round $\ell-1$, each active node adds Laplace noise scaled to $\|\mathcal{X}^{(\ell-1)}\|_\infty$.
A key distinction from walk counting is that, under random marking, each edge affects at most one node’s computation in a single round. 
Hence, by parallel composition (Lemma~\ref{lem:pp}), the privacy budget $\varepsilon$ does not need to be divided across nodes or rounds. 



\begin{figure*}[htbp!]
  \centering
  \includegraphics[width=0.96\linewidth]{figures/3-lp.pdf} % file myfig.svg in figs/
  \caption{Illustration of the $3$-line path counting mechanism $\mathcal{M}_{\mathrm{3\text{-}pt}}$ applied to graph $G$ with privacy budget $\varepsilon =1$.
For simplicity, $\sum$ denotes $\sum_{j \in \mathcal{N}(i)}$.
The initialization step occurs before random marking, but is shown afterward for clarity. 
}
\label{fig:3-lp}
\end{figure*}



With the above idea, we establish the $k$-round edge-LDP mechanism $\mathcal{M}_{\mathrm{k\text{-}pt}}$ for $k$-line path counting. 
Compared with the non-private local client $\widehat{\mathcal{R}}^{(\ell)}$ in Algorithm~\ref{alg:k-lp_nondp}, here the local randomizer $\mathcal{R}^{(\ell)}$ adds noise $\mathrm{Lap}(\|\mathcal{X}^{(\ell-1)}\|_\infty/\varepsilon)$ in Line~9 and noise $\mathrm{Lap}(1/\varepsilon)$ to the multiplicative factor 
$\sum_{j \in \mathcal{N}(i)} \mathbf{1}\{r_j = k\}$ in Line~11. 
{When evaluating the privacy of $\mathcal{M}_{\mathrm{k\text{-}pt}}$, astute readers may consider privacy amplification, whereby privacy guarantees can be strengthened when the DP result is derived from an unknown subsample of the input dataset.
However, privacy amplification does not apply here. 
In $\mathcal{M}_{\mathrm{k\text{-}pt}}$, node markings are publicly revealed, making it deterministically known whether an edge will be sampled and thus violating the sampling uncertainty required for privacy amplification.
Moreover, privatizing node markings to enable privacy amplification introduces new challenges, as each mark must be revealed to neighboring nodes during communication.
As a result, we establish the following privacy guarantee:
}
\begin{theorem}
\label{the:k-lp_privacy}
    For any graph $G$, path length $k$, and privacy budget $\varepsilon$, the mechanism $\mathcal{M}_{\mathrm{k\text{-}pt}}$ satisfies $\varepsilon$-edge-LDP.
\end{theorem}
For utility analysis, the total error $|\mathcal{M}_{\mathrm{k\text{-}pt}}(G, \varepsilon)- P_k|$ consists of two parts.
The first is the sampling error incurred by random marking, as shown in Lemma~\ref{lem:k-lp_sampling_bound}.
For the second part, the {DP noise} introduced by the Laplace mechanism, we show that it slightly differs from that of $k$-line walk counting but admits the same asymptotic upper bound.
\begin{lemma}
    \label{lem:k-lp_additive_bound}
    For any graph $G$, path length $k$, parameters $\varepsilon$ and $\beta$, 
$\mathcal{M}_{\mathrm{k\text{-}pt}}$ satisfies
$
    \mathbb{E}[\mathcal{M}_{\mathrm{k\text{-}pt}}(G, \varepsilon)]
    =  
    \widehat{\mathcal{M}}_{\mathrm{k\text{-}pt}}(G)$
under the same node marking $(r_1, \dots, r_N)$, 
and with probability at least $1 - \beta$,
\begin{align*}
    |\mathcal{M}_{\mathrm{k\text{-}pt}}(G, \varepsilon) - 
    \widehat{\mathcal{M}}_{\mathrm{k\text{-}pt}}(G)| 
    &\leq 
    {k(k+1)} \,
    \gamma
    \sqrt{N}
    (\hat{d}(G) + \gamma \sqrt{\hat{d}(G)} + \gamma)^{k-1} 
    \\
    &=  \tilde{O}(\sqrt{N} d(G)^{k-1}),
\end{align*}
where $\gamma = (k+1)\sqrt{8\log(6N/\beta)}/\varepsilon = \tilde{O}(1) $ and $\hat{d}(G) = \max(d(G), \lceil\gamma^2\rceil)$.
\end{lemma}
% \orange{Accounting for both the sampling error from random marking and the DP noise, we obtain the following utility bound:}
Combining both parts of the error, we establish the utility bound:
\begin{theorem}
\label{the:k-lp_bound}
    For any graph $G$, path length $k$, parameters $\varepsilon$ and $\beta$, the $k$-round mechanism $\mathcal{M}_{\mathrm{k\text{-}pt}}$ satisfies $
    \mathbb{E}[\mathcal{M}_{\mathrm{k\text{-}pt}}(G, \varepsilon)] =  
    P_k$,
and with probability at least $1 - \beta$,
\begin{align*}
    |\mathcal{M}_{\mathrm{k\text{-}pt}}(G, \varepsilon) - 
    P_k|  \leq  (k+1) \sqrt{2N\beta^{-1}}(d(G)+k+1)^k + \\ 
          {k(k+1)} \, 
        \gamma
        \sqrt{N}
        (\hat{d}(G) + \gamma \sqrt{\hat{d}(G)} + \gamma)^{k-1} 
          =  \tilde{O}(\sqrt{N} d(G)^{k}),
\end{align*}
where $\gamma = (k+1)\sqrt{8\log(12N/\beta)}/\varepsilon = \tilde{O}(1) $ and $\hat{d}(G) = \max(d(G), \lceil\gamma^2\rceil)$.
\end{theorem}



In summary, $\mathcal{M}_{\mathrm{k\text{-}pt}}$ still produces an unbiased estimator of the number of $k$-line paths $P_k$, while incurring a $\tilde{O}(d(G))$ amplification in error upper bound compared to the $k$-line walk counting solution.
Specifically, the overall error is asymptotically dominated by the sampling error, while the {DP noise} matches that of walk counting. 




\paragraph{{Remark}}
{
Notably, unbiased estimators are common in LDP settings, especially for RR-based mechanisms, where the final result can be regarded as a sum of independent random variables, and unbiasedness follows directly.
However, our mechanism induces correlations among different noise sources, which significantly complicates proving unbiasedness.}

{
Besides, in our mechanism, the sampling error also arises from the privacy aspect, even though it is independent of the privacy budget $\varepsilon$. 
This is because random marking is not specific to path counting; instead, it addresses an LDP-specific challenge, i.e., eliminating duplicate nodes when each node observes only limited local information.
While our theoretical analysis suggests that the sampling error asymptotically dominates the DP noise by a factor of $\tilde{O}(d(G))$, our experiments show that, in practice, the gap is much smaller. 
{
This observation indicates that applying random marking offers a near ``free lunch'' in our mechanism. 
}
Moreover, we can optimize the overall error by mitigating this gap.
Since both mechanisms $\widehat{\mathcal{M}}_{\mathrm{k\text{-}pt}}$ and ${\mathcal{M}}_{\mathrm{k\text{-}pt}}$ produce unbiased estimators, we run the LDP mechanism
$n_{\text{rep}}$ times with per run privacy budget $\varepsilon/n_{\text{rep}}$ and average the outputs.
This reduces the sampling error, and while allocating a smaller per-run privacy budget increases the DP noise, averaging across runs mitigates its impact. 
We demonstrate the effectiveness of this optimization strategy in
Section~\ref{subsec:sample_dp}.
}




% \paragraph{\red{Remark}}
% \marginnote{\red{R1.O1}}
% \red{
% Importantly, the sampling error should be regarded as a component of the LDP error despite its independence from the privacy budget $\varepsilon$.
% More precisely, this is because the random marking technique is not designed specifically for the path counting; rather, it is used to address a unique challenge in the LDP setting, i.e., eliminating duplicate nodes in the collected paths when each node has access only to limited local information. 
% Although the sampling error appears to asymptotically dominate the DP noise with a theoretical gap of $O(d(G))$, empirical evaluation results in our experiment show that these two parts are much smaller than $d(G)$.
% Furthermore, we can optimize the overall error by mitigating this gap: 
% Since both mechanisms $\widehat{\mathcal{M}}_{\mathrm{k\text{-}pt}}$ and ${\mathcal{M}}_{\mathrm{k\text{-}pt}}$ produce unbiased estimators, we running the LDP mechanism $n_{\text{rep}}$ times, each with a privacy budget of $\varepsilon / n_{\text{rep}}$, and averaging the $n_{\text{rep}}$ outputs to obtain the final estimate.
% Since the sampling error is independent of $\varepsilon$, this strategy effectively reduces it through repeated execution. 
% Although allocating a smaller privacy budget to each run increases the DP error, averaging the outputs across multiple runs helps mitigate this increase.
% We will show the effectiveness of this strategy in the experiment in Section~\ref{subsec:sample_dp}. 
% }







\paragraph{Optimization in efficiency}

Earlier in this section, we required that nodes in rounds inconsistent with their marks report $0$. 
In fact, this step can be omitted, since the computation results depend only on non-zero values. 
As a result, we instead require that nodes remain silent in such rounds.
Moreover, since only nodes with mark $\ell$ are active in round $\ell$, communication can be optimized by having the analyzer $\mathcal{A}$ send $\|\mathcal{X}^{(\ell-1)}\|_\infty$ exclusively to those nodes; their outputs are then forwarded only to neighbors whose mark is $\ell+1$.
This requires the analyzer to know each node’s mark, which can be achieved by modifying Line~6 of Algorithm~\ref{alg:k-lp_nondp} so that nodes also send their marks to $\mathcal{A}$. 
% Both optimizations jointly achieve a constant-factor reduction in communication between nodes and with the analyzer. 



Figure~\ref{fig:3-lp} shows an example of executing the optimized 3-round mechanism $\mathcal{M}_{\mathrm{3\text{-}pt}}$ on graph $G$ with $\varepsilon = 1$, showing only the active nodes at each step.
The mechanism first initializes the number of $0$-line paths ending at each node to $1$.
After that, each node samples its mark and becomes active only in the corresponding round. 
Since initialization corresponds to round $0$,  
only outputs of nodes with mark $0$ are used; outputs from all other nodes are unused and shown in gray.
% from all other nodes are never used and are shown grayed out.
In the final round, the analyzer aggregates the outputs of all active nodes and rescales the sum to produce the final result $19.2$.






\paragraph{Communication and computation analysis}  
The total communication cost of $\mathcal{M}_{\mathrm{k\text{-}pt}}$ is $O(M+N)$. Communication between nodes contributes $O(M)$, as each node sends its mark and noisy result once to all its neighbors. 
The communication between nodes and the analyzer adds an additional $O(N)$, as each node sends its mark and participates in at most one round of receiving and sending values with the analyzer.
In terms of computational cost, the analyzer runs in $O(N)$ time across $k-2$ rounds of computing the round maximum and one round of aggregating the final result. 
On the node side, the runtime is $O(d(G))$, as each node participates in one round of aggregating $O(d(G))$ values from its neighbors.





